\hypertarget{struttura-di-applicazioni-complesse}{%
\section{Struttura di applicazioni
complesse}\label{struttura-di-applicazioni-complesse}}

Nelle applicazioni viste finora, la gestione di \emph{cosa fare} in base
al differente tipo di richiesta (GET, POST) e dei parametri ricevuti,
veniva gestita da un unico file, contenente una serie di \texttt{if},
ma, anche con poche e semplici funzionalità, tendeva a intricarsi
velocemente.

Per gestire meglio l'organizzazione di applicazioni di media
complessità, senza aver la pretesa di utilizzare librerie complete ma
anche decisamente più complesse, si vedrà come utilizzare Slim e se ne
fornirà una proposta di come strutturare applicazioni con accesso a
database e con presenza di più funzionalità.

\hypertarget{cosuxe9-slim}{%
\subsection{Cos'é Slim}\label{cosuxe9-slim}}

Il sito di Slim definisce la libreria in questo modo

\begin{quote}
Slim is a PHP micro framework that helps you quickly write simple yet
powerful web applications and APIs. At its core, Slim is a dispatcher
that receives an HTTP request, invokes an appropriate callback routine
and returns an HTTP response. That's it.
\end{quote}

Quindi il nocciolo di quello che fa è di ricevere le richieste HTTP,
\emph{agganciarle} all'azione appropriata e ritornare la corretta
risposta HTTP.

\hypertarget{come-creare-la-prima-applicazione-in-slim}{%
\subsection{Come creare la prima applicazione in
Slim}\label{come-creare-la-prima-applicazione-in-slim}}

Per informazioni più dettagliate si veda il sito di riferimento:
\href{https://www.slimframework.com/docs/v4/}{Slim}

\hypertarget{come-installare}{%
\subsubsection{Come installare}\label{come-installare}}

Utilizzando Composer\footnote{Nei laboratori dove è presente il proxy, è
  necessario operare con Fiddler come si è fatto in altre occasioni}
bisogna installare nel proprio progetto la libreria \texttt{Slim}

\begin{Shaded}
\begin{Highlighting}[]
\ExtensionTok{composer}\NormalTok{ require slim/slim:}\StringTok{"4.*"}
\end{Highlighting}
\end{Shaded}

Successivamente è necessario un altro componente per gestire le
richieste, che va installato sempre con Composer.

\begin{Shaded}
\begin{Highlighting}[]
\ExtensionTok{composer}\NormalTok{ require slim/psr7}
\end{Highlighting}
\end{Shaded}

A questo punto creare nel progetto una cartella \emph{public} e al suo
interno inserire il file \emph{index.php}, che diventerà il
front-controller dell'applicazione, cioè ogni richiesta passerà
attraverso di lui.

Il contenuto di questo file può essere popolato in questo modo

\begin{Shaded}
\begin{Highlighting}[]
\KeywordTok{\textless{}?php}
\KeywordTok{use} \ConstantTok{P}\NormalTok{sr\textbackslash{}}\ConstantTok{H}\NormalTok{ttp\textbackslash{}}\ConstantTok{M}\NormalTok{essage\textbackslash{}}\ConstantTok{R}\NormalTok{esponseInterface }\KeywordTok{as} \ConstantTok{R}\NormalTok{esponse}\OtherTok{;}
\KeywordTok{use} \ConstantTok{P}\NormalTok{sr\textbackslash{}}\ConstantTok{H}\NormalTok{ttp\textbackslash{}}\ConstantTok{M}\NormalTok{essage\textbackslash{}}\ConstantTok{S}\NormalTok{erverRequestInterface }\KeywordTok{as} \ConstantTok{R}\NormalTok{equest}\OtherTok{;}
\KeywordTok{use} \ConstantTok{S}\NormalTok{lim\textbackslash{}}\ConstantTok{F}\NormalTok{actory\textbackslash{}}\ConstantTok{A}\NormalTok{ppFactory}\OtherTok{;}

\KeywordTok{require} \ConstantTok{\_\_DIR\_\_} \OperatorTok{.} \StringTok{\textquotesingle{}/../vendor/autoload.php\textquotesingle{}}\OtherTok{;}

\VariableTok{$app} \OperatorTok{=} \ConstantTok{A}\NormalTok{ppFactory::create()}\OtherTok{;}

\CommentTok{//Questa parte deve contenere il percorso della}
\CommentTok{//sottocartella dove si trova l\textquotesingle{}applicazione in questo caso inserito nella}
\CommentTok{//variabile di configurazione BASE\_PATH}
\VariableTok{$app}\NormalTok{{-}\textgreater{}setBasePath(}\ConstantTok{BASE\_PATH}\NormalTok{)}\OtherTok{;}

\VariableTok{$app}\NormalTok{{-}\textgreater{}get(}\StringTok{\textquotesingle{}/\textquotesingle{}}\OtherTok{,} \KeywordTok{function}\NormalTok{ (}\ConstantTok{R}\NormalTok{equest }\VariableTok{$request}\OtherTok{,} \ConstantTok{R}\NormalTok{esponse }\VariableTok{$response}\OtherTok{,} \VariableTok{$args}\NormalTok{) \{}
    \VariableTok{$response}\NormalTok{{-}\textgreater{}getBody(){-}\textgreater{}write(}\StringTok{\textquotesingle{}Hello world!\textquotesingle{}}\NormalTok{)}\OtherTok{;}
    \ControlFlowTok{return} \VariableTok{$response}\OtherTok{;}
\NormalTok{\})}\OtherTok{;}

\VariableTok{$app}\NormalTok{{-}\textgreater{}get(}\StringTok{\textquotesingle{}/altra\_pagina\textquotesingle{}}\OtherTok{,} \KeywordTok{function}\NormalTok{ (}\ConstantTok{R}\NormalTok{equest }\VariableTok{$request}\OtherTok{,} \ConstantTok{R}\NormalTok{esponse }\VariableTok{$response}\OtherTok{,} \VariableTok{$args}\NormalTok{) \{}
    \VariableTok{$response}\NormalTok{{-}\textgreater{}getBody(){-}\textgreater{}write(}\StringTok{\textquotesingle{}Questa è un}\SpecialCharTok{\textbackslash{}\textquotesingle{}}\StringTok{altra pagina\textquotesingle{}}\NormalTok{)}\OtherTok{;}
    \ControlFlowTok{return} \VariableTok{$response}\OtherTok{;}
\NormalTok{\})}\OtherTok{;}

\VariableTok{$app}\NormalTok{{-}\textgreater{}run()}\OtherTok{;}
\end{Highlighting}
\end{Shaded}

Le parti più importanti di questo codice sono: 1. la variabile
\texttt{\$app} rappresenta l'applicazione ed è quella che verrà usata
per creare tutte le associazioni rotta - azione 2. per ogni rotta che si
intende gestire, che in questo esempio sono la rotta ``vuota''
\texttt{/} e la rotta \texttt{/altra\_pagina}, deve essere chiamato un
metodo che mappa la richiesta HTTP che si intende utilizzare (GET, POST,
\ldots) e nel quale vanno passati come parametri la rotta e l'azione da
associare, sotto forma di funzione anonima. In questo esempio viene
utilizzato per entrambe il metodo \texttt{get} e quindi una GET sulla
rotta vuota produrrà una pagina con la scritta \texttt{Hello\ world} e
una GET sulla rotta \texttt{/altra\_pagina} produrrà la scritta
\texttt{Questa\ è\ un\textquotesingle{}altra\ pagina}. 3. I parametri
della funzione anonima sono tre e rappresentano nell'ordine l'oggetto
che contiene le informazioni della richiesta (\texttt{\$request}),
quello che conterrà le informazioni della risposta (\texttt{\$response})
e un vettore (\texttt{\$args}) che conterrà eventuali parametri
contenuti nella rotta (si veda l'esempio con i template) 4. Nel corpo
della funzione avviene la scrittura del body della risposta, utilizzando
il metodo \texttt{write} applicato appunto al body e scrivendo quello
che si vuole sia contenuto nella pagina di risposta
(\textbf{attenzione}: il body della risposta è quello relativo al
protocollo HTTP, non ha nulla a che vedere con il body della pagina
HTML, hanno solo lo stesso nome).

\textbf{Attenzione}: il parametro nella chiamata al metodo
\texttt{setBasePath} deve rispecchiare l'esatta struttura delle cartelle
dove si trova l'applicazione. In questo esempio il metodo viene chiamato
in questo modo

\begin{Shaded}
\begin{Highlighting}[]
\VariableTok{$app}\NormalTok{{-}\textgreater{}setBasePath(}\StringTok{"/template\_slim"}\NormalTok{)}\OtherTok{;}
\end{Highlighting}
\end{Shaded}

perché l'applicazione si trova nella sottocartella \emph{template\_slim}
all'interno della root del web server. Se ad esempio la cartella si
trovasse nella cartella \emph{\textasciitilde piero/negozio} allora
l'istruzione andrebbe riscritta come

\begin{Shaded}
\begin{Highlighting}[]
\VariableTok{$app}\NormalTok{{-}\textgreater{}setBasePath(}\StringTok{"/\textasciitilde{}piero/negozio"}\NormalTok{)}\OtherTok{;}
\end{Highlighting}
\end{Shaded}

\hypertarget{configurare-apache-per-lurl-rewriting}{%
\subsection{Configurare Apache per l'URL
rewriting}\label{configurare-apache-per-lurl-rewriting}}

In un'applicazione web è una buona idea avere un solo
\emph{front-controller}, cioè un unico file, nel nostro caso
\texttt{index.php}, a cui arrivano tutte le richieste HTTP,
indipendentemente dalla \emph{rotta} effettiva della richiesta: quindi
richieste indirizzate ad esempio a \texttt{/negozio/accessori} piuttosto
che ad \texttt{/admin/} passeranno comunque attraverso il file
\texttt{index.php}, che si occuperà del \emph{dispatching}, della
\emph{distribuzione}, al pezzo di codice che dovrà gestire ogni
particolare richiesta.

Se quindi l'applicazione ha bisogno che tutte le richieste HTTP arrivino
al file \emph{index.php}, è necessario istruire il web server per fare
in modo che ogni richiesta che arriva, indipendentemente dall'URL
effettivo, venga dirottata sul file \texttt{index.php}.

Usando Apache presente in Xampp bisogna seguire questi passi: -
abilitare il modulo di URL rewriting. Per far questo aprire con un
editor di testo (anche PHPStorm va bene) il file
\texttt{directory\_di\_xampp/apache/conf/httpd.conf} e controllare se la
riga

\begin{Shaded}
\begin{Highlighting}[]
\NormalTok{LoadModule rewrite\_module modules/mod\_rewrite.so}
\end{Highlighting}
\end{Shaded}

ha un \textbf{\#} davanti. Se così fosse va eliminato e altrettanto va
fatto per la riga

\begin{Shaded}
\begin{Highlighting}[]
\NormalTok{LoadModule actions\_module modules/mod\_actions.so}
\end{Highlighting}
\end{Shaded}

A questo punto bisogna far partire o ripartire Apache e la
configurazione dovrebbe essere corretta (per ulteriori informazioni
consultare
\href{https://www.slimframework.com/docs/v4/start/web-servers.html}{questa
pagina} oppure Google).

Infine, per fare in modo che tutte le richieste vengano dirottate sul
file \texttt{index.php} che si trova nella cartella
\texttt{public}\footnote{Il motivo per il quale il file
  \texttt{index.php} viene inserito nella cartella \texttt{public} e non
  nella root del progetto è legato alla sicurezza dell'applicazione,
  anche se qui non verrà approfondito questo tema.} bisogna aggiungere
nella cartella principale, \emph{template\_slim} in questo esempio, un
file \emph{.htaccess} con il seguente contenuto:

\begin{Shaded}
\begin{Highlighting}[]
\NormalTok{RewriteEngine On}

\NormalTok{RewriteRule \^{} public/index.php [QSA,L]}
\end{Highlighting}
\end{Shaded}

Adesso dovrebbe essere sufficiente inserire l'indirizzo
\texttt{localhost/template\_slim} (o quello relativo a dove si è deciso
di mettere l'applicazione) nella barra degli indirizzi del browser e
dovrebbe apparire una pagina con la scritta \textbf{Hello world}. Per
ulteriore conferma provare a inserire
\texttt{localhost/slim/altra\_pagina} e dovrebbe comparire il testo
\textbf{Questa è un'altra pagina}.

\hypertarget{gestione-degli-errori}{%
\subsection{Gestione degli errori}\label{gestione-degli-errori}}

Slim prevede già un meccanismo di gestione degli errori/eccezioni con un
proprio Middleware\footnote{Per \emph{middleware} in questi contesti si
  intende uno strato di software che si interpone tra la richiesta
  dell'utente e l'applicazione vera e proprio, per svolgere delle
  funzionalità che sono comuni a moltissime applicazioni e che, quindi,
  possono essere riutilizzate senza bisogno di particolari modifiche.
  Esempi di middleware sono appunto la gestione delle eccezioni e
  l'autenticazione.} e per attivarlo è sufficiente introdurre questa
riga di codice

\begin{Shaded}
\begin{Highlighting}[]
\VariableTok{$errorMiddleware} \OperatorTok{=} \VariableTok{$app}\NormalTok{{-}\textgreater{}addErrorMiddleware(}\KeywordTok{true}\OtherTok{,} \KeywordTok{true}\OtherTok{,} \KeywordTok{true}\NormalTok{)}\OtherTok{;}
\end{Highlighting}
\end{Shaded}

Di default il comportamento è quello di intercettare errori/eccezioni e
mandare in risposta una pagina HTML con le informazioni sul problema,
eventualmente con anche i dettagli dello stack di chiamate nel caso il
primo parametro fosse uguale a \texttt{true}, come nell'esempio qua
sopra. Gli altri due parametri sono rilevanti solo se si utilizza un
meccanismo di \emph{logging}, per maggiori informazioni fare riferimento
alla
\href{https://www.slimframework.com/docs/v4/middleware/error-handling.html}{documentazione
ufficiale}.

\hypertarget{utilizzare-plates}{%
\subsection{Utilizzare Plates}\label{utilizzare-plates}}

Volendo utilizzare la libreria di templates Plates, come già fatto in
precedenza, dovrà essere importata con il solito comando

\begin{Shaded}
\begin{Highlighting}[]
\ExtensionTok{composer}\NormalTok{ require league/plates}
\end{Highlighting}
\end{Shaded}

e nel codice andrà indicata con

\begin{Shaded}
\begin{Highlighting}[]
\KeywordTok{use} \ConstantTok{L}\NormalTok{eague\textbackslash{}}\ConstantTok{P}\NormalTok{lates\textbackslash{}}\ConstantTok{E}\NormalTok{ngine}\OtherTok{;}
\end{Highlighting}
\end{Shaded}

ed è poi buona cosa utilizzare un
\href{https://en.wikipedia.org/wiki/Dependency_injection}{Dipendency
Injector} per utilizzarla all'interno di Slim. Senza dilungarsi sul
concetto di Dipendency Injection, lo scopo è quello di separare in
maniera pulita un servizio (in questo caso Plates) dall'applicazione
vera e propria che utilizzerà quel servizio. In pratica per ottenere lo
scopo bisogna prima importare una libreria che implementi il meccanismo
di DI, e Slim suggerisce di utilizzare PHP-DI, e con Composer si fa

\begin{Shaded}
\begin{Highlighting}[]
\ExtensionTok{composer}\NormalTok{ require php{-}di/slim{-}bridge:}\PreprocessorTok{*} \AttributeTok{{-}{-}with{-}all{-}dependencies}
\end{Highlighting}
\end{Shaded}

Fatto questo, basterà aggiungere nel codice la creazione del Container,
che poi dovrà essere passato all'applicazione

\begin{Shaded}
\begin{Highlighting}[]
\VariableTok{$container} \OperatorTok{=} \KeywordTok{new} \ConstantTok{C}\NormalTok{ontainer()}\OtherTok{;}
\CommentTok{//da inserire prima della create di AppFactory}
\ConstantTok{A}\NormalTok{ppFactory::setContainer(}\VariableTok{$container}\NormalTok{)}\OtherTok{;}
\end{Highlighting}
\end{Shaded}

Per aggiungere dei servizi al container è sufficiente usare il metodo
\texttt{set}, che nel caso specifico di Plates verrà fatto in questo
modo

\begin{Shaded}
\begin{Highlighting}[]
\VariableTok{$container}\NormalTok{{-}\textgreater{}set(}\StringTok{\textquotesingle{}template\textquotesingle{}}\OtherTok{,} \KeywordTok{function}\NormalTok{ ()\{}
    \VariableTok{$engine} \OperatorTok{=} \KeywordTok{new} \ConstantTok{E}\NormalTok{ngine(}\StringTok{\textquotesingle{}../templates\textquotesingle{}}\OtherTok{,} \StringTok{\textquotesingle{}tpl\textquotesingle{}}\NormalTok{)}\OtherTok{;}
    \VariableTok{$engine}\NormalTok{{-}\textgreater{}addData([}\StringTok{\textquotesingle{}base\_path\textquotesingle{}}\NormalTok{ =\textgreater{} }\ConstantTok{BASE\_PATH}\NormalTok{])}\OtherTok{;}
    \ControlFlowTok{return} \VariableTok{$engine}\OtherTok{;}
\NormalTok{\})}\OtherTok{;}
\end{Highlighting}
\end{Shaded}

dove il primo parametro è la chiave per recuperare il servizio e il
secondo è la funzione anonima che ritorna il servizio, occupandosi della
sua creazione.

\textbf{Attenzione}: il primo parametro del costruttore dell'Engine deve
come al solito contenere la cartella dove poi verranno inseriti i vari
template, ma il percorso deve essere relativo al file dove si trovano
queste istruzioni: nel nostro caso in cui il file \texttt{index.php} si
trova nella cartella \textbf{public} il nome della cartella deve essere
preceduto da \texttt{../}. Inoltre viene usato anche il metodo
\texttt{add} per aggiungere una variabile, in questo caso
\texttt{base\_path}, in modo che sia visibile al template. Questa
variabile è necessaria per poter generare dei link corretti all'interno
dei template e contiene

Una volta aggiunto Plates nel container, per poterlo utilizzare basterà
aggiungere del codice come il seguente nelle callback per le rotte su
cui vogliamo usare dei template.

\begin{Shaded}
\begin{Highlighting}[]
\VariableTok{$app}\NormalTok{{-}\textgreater{}get(}\StringTok{\textquotesingle{}/saluti/\{name\}\textquotesingle{}}\OtherTok{,} \KeywordTok{function}\NormalTok{ (}\ConstantTok{R}\NormalTok{equest }\VariableTok{$request}\OtherTok{,} \ConstantTok{R}\NormalTok{esponse }\VariableTok{$response}\OtherTok{,} \VariableTok{$args}\NormalTok{) \{}
    \VariableTok{$template} \OperatorTok{=} \VariableTok{$this}\NormalTok{{-}\textgreater{}get(}\StringTok{\textquotesingle{}template\textquotesingle{}}\NormalTok{)}\OtherTok{;}
    \VariableTok{$response}\NormalTok{{-}\textgreater{}getBody(){-}\textgreater{}write(}\VariableTok{$template}\NormalTok{{-}\textgreater{}render(}\StringTok{\textquotesingle{}saluti\textquotesingle{}}\OtherTok{,}
\NormalTok{            [}
                \StringTok{\textquotesingle{}name\textquotesingle{}}\NormalTok{ =\textgreater{} }\VariableTok{$args}\NormalTok{[}\StringTok{\textquotesingle{}name\textquotesingle{}}\NormalTok{]}
\NormalTok{            ]}
\NormalTok{        )}
\NormalTok{    )}\OtherTok{;}
    \ControlFlowTok{return} \VariableTok{$response}\OtherTok{;}
\NormalTok{\})}\OtherTok{;}
\end{Highlighting}
\end{Shaded}

In questo esempio ci sono due aspetti della libreria \texttt{Slim} che
devono essere approfonditi: 1. La rotta contiene al suo interno una
parte parametrica, inclusa tra le parentesi graffe, in questo esempio
\emph{name}. Quel parametro indica che una rotta in cui la prima parte è
\texttt{/saluti/} verrà intercettata. 2. La parte parametrica verrà
recuperata attraverso il parametro \texttt{\$args}, che sarà la libreria
stessa a popolare con gli eventuali argomenti passati attraverso la GET,
in questo caso appunto il \texttt{name}. Se quindi la rotta fosse
\texttt{GET\ /saluti/alessandro}, il vettore \texttt{\$args} avrà una
chiave di nome \texttt{name} il cui valore sarà \texttt{alessandro},
valore che potrà essere recuperato come si può vedere nel codice con la
solita notazione
\texttt{\$args{[}\textquotesingle{}name\textquotesingle{}{]}}.

\hypertarget{connessione-al-database}{%
\subsection{Connessione al database}\label{connessione-al-database}}

Per utilizzare il database e dare un minimo di organizzazione, conviene
isolare l'aspetto della connessione in una classe apposita, che avrà un
metodo da invocare ogni volta che nell'applicazione è necessario
recuperare dei dati dal DB. Per far questo vien creato un file,
\emph{config.php} all'interno della cartella \emph{conf}, che conterrà
le informazioni necessarie alla connessione in questo modo:

\begin{Shaded}
\begin{Highlighting}[]
\KeywordTok{const} \ConstantTok{DB\_HOST} \OperatorTok{=} \StringTok{\textquotesingle{}host\textquotesingle{}}\OtherTok{;}
\KeywordTok{const} \ConstantTok{DB\_NAME} \OperatorTok{=} \StringTok{\textquotesingle{}db\_name\textquotesingle{}}\OtherTok{;}
\KeywordTok{const} \ConstantTok{DB\_USER} \OperatorTok{=} \StringTok{\textquotesingle{}user\textquotesingle{}}\OtherTok{;}
\KeywordTok{const} \ConstantTok{DB\_PASSWORD} \OperatorTok{=} \StringTok{\textquotesingle{}password\textquotesingle{}}\OtherTok{;}
\KeywordTok{const} \ConstantTok{DB\_CHAR} \OperatorTok{=} \StringTok{\textquotesingle{}utf8\textquotesingle{}}\OtherTok{;}
\end{Highlighting}
\end{Shaded}

con i parametri opportunamente configurati per la propria connessione.

La classe che fornirà la connessione vera e propria verrà invece creata
nel file \emph{Connection.php} all'interno della cartella \emph{Util}.

Per poter utilizzare in ogni punto dell'applicazione questa nuova classe
sarà necessario indicare al meccanismo di \emph{autoloading} dove
recuperarla, in modo da non sia necessario dover includere a mano i file
ogni volta. Per far questo nel file \emph{Connection.php} sarà
necessario aggiungere la linea

\begin{Shaded}
\begin{Highlighting}[]
\KeywordTok{namespace} \ConstantTok{U}\NormalTok{til}\OtherTok{;}
\end{Highlighting}
\end{Shaded}

e poi ricostruire i file di autoloading aggiungendo le righe

\begin{Shaded}
\begin{Highlighting}[]
  \StringTok{"autoload"}\OtherTok{:}\NormalTok{ \{}
    \StringTok{"psr{-}4"}\OtherTok{:}\NormalTok{ \{}
      \StringTok{"Util}\SpecialCharTok{\textbackslash{}\textbackslash{}}\StringTok{"} \OtherTok{:} \StringTok{"Util/"}
\NormalTok{    \}}
\NormalTok{  \}}\OtherTok{,}
\end{Highlighting}
\end{Shaded}

all'interno di \emph{composer.json} e chiamando \textbf{composer} in
questo modo

\begin{Shaded}
\begin{Highlighting}[]
\ExtensionTok{composer}\NormalTok{ dumpautoload}
\end{Highlighting}
\end{Shaded}

in modo che aggiorni i file all'interno della cartella \emph{vendor} per
permettere il caricamento della classe \emph{Connection}, che quindi
sarà utilizzabile semplicemente aggiungendo la riga

\begin{Shaded}
\begin{Highlighting}[]
\KeywordTok{use} \ConstantTok{U}\NormalTok{til\textbackslash{}}\ConstantTok{C}\NormalTok{onnection}\OtherTok{;}
\end{Highlighting}
\end{Shaded}

all'inizio di ogni file che la vorrà utilizzare, nel nostro caso
\emph{index.php}.

La classe \emph{Connection} avrà un solo attributo statico,
\texttt{\$pdo}, che rappresenta appunto la connessione al DB, avrà il
costruttore privato perchè non è pensata per essere istanziata e un
metodo \texttt{getInstance()} che restituisce la connessione,
istanziando la connessione se è la prima volta che viene chiamato in uno
script oppure restituendo la connessione corrente se esiste già. Si veda
il contenuto di \emph{Connection.php} per maggiori dettagli.

\hypertarget{gestione-dei-dati-presenti-nel-db-con-il-pattern-mvc}{%
\subsection{Gestione dei dati presenti nel DB con il pattern
MVC}\label{gestione-dei-dati-presenti-nel-db-con-il-pattern-mvc}}

Dove la View in queste applicazioni è il \emph{template}, per quanto
riguarda il Model si procederà come visto in precedenza, dove per il
controller si creeranno delle classi in base alle esigenze
dell'applicazione.

\hypertarget{esempio-della-gestione-di-un-negozio}{%
\subsubsection{Esempio della gestione di un
negozio}\label{esempio-della-gestione-di-un-negozio}}

Per applicare i concetti a un esempio verrà ripreso il negozio di
prodotti d'abbigliamento, dove per semplicità verrà considerata solo la
presenza della tabella \textbf{prodotto} mostrata qua sotto:

\begin{Shaded}
\begin{Highlighting}[]
\NormalTok{erDiagram}
\NormalTok{    prodotto \{}
\NormalTok{string id PK}
\NormalTok{string nome}
\NormalTok{string descrizione}
\NormalTok{float prezzo}
\NormalTok{string genere "Uomo o donna"}
\NormalTok{\}}
\end{Highlighting}
\end{Shaded}

\hypertarget{la-classe-prodottorepository}{%
\subsection{La classe
ProdottoRepository}\label{la-classe-prodottorepository}}

L'implementazione del modello verrà realizzata creando una classe (o più
classi) apposita che contenga al proprio interno i metodi ritenuti
opportuni per l'interazione con i dati. Questi metodi saranno
\emph{statici}, in quanto la classe non verrà utilizzata per creare
oggetti, ma solo per avere un'organizzazione dei metodi ordinata.

Per la tabella \texttt{prodotto} un possibile insieme di metodi, con la
relativa implementazione, potrebbe essere il seguente:

\begin{Shaded}
\begin{Highlighting}[]
\KeywordTok{class} \ConstantTok{P}\NormalTok{rodottoRepository\{}

    \KeywordTok{public} \KeywordTok{static} \KeywordTok{function}\NormalTok{ listAll()\{}
        \VariableTok{$pdo} \OperatorTok{=} \ConstantTok{C}\NormalTok{onnection::getInstance()}\OtherTok{;}
        \VariableTok{$risposta} \OperatorTok{=} \VariableTok{$pdo}\NormalTok{{-}\textgreater{}query(}\StringTok{"SELECT * FROM prodotto"}\NormalTok{)}\OtherTok{;}
        \ControlFlowTok{return} \VariableTok{$risposta}\NormalTok{{-}\textgreater{}fetchAll()}\OtherTok{;}
\NormalTok{    \}}

    \KeywordTok{public} \KeywordTok{static} \KeywordTok{function}\NormalTok{ listAllMale()\{}
        \VariableTok{$pdo} \OperatorTok{=} \ConstantTok{C}\NormalTok{onnection::getInstance()}\OtherTok{;}
        \VariableTok{$risposta} \OperatorTok{=} \VariableTok{$pdo}\NormalTok{{-}\textgreater{}query(}\StringTok{\textquotesingle{}SELECT * FROM prodotto WHERE genere = "Uomo"\textquotesingle{}}\NormalTok{)}\OtherTok{;}
        \ControlFlowTok{return} \VariableTok{$risposta}\NormalTok{{-}\textgreater{}fetchAll()}\OtherTok{;}
\NormalTok{    \}}

    \KeywordTok{public} \KeywordTok{static} \KeywordTok{function}\NormalTok{ listAllFemale()\{}
        \VariableTok{$pdo} \OperatorTok{=} \ConstantTok{C}\NormalTok{onnection::getInstance()}\OtherTok{;}
        \VariableTok{$risposta} \OperatorTok{=} \VariableTok{$pdo}\NormalTok{{-}\textgreater{}query(}\StringTok{\textquotesingle{}SELECT * FROM prodotto WHERE genere = "Donna"\textquotesingle{}}\NormalTok{)}\OtherTok{;}
        \ControlFlowTok{return} \VariableTok{$risposta}\NormalTok{{-}\textgreater{}fetchAll()}\OtherTok{;}
\NormalTok{    \}}
\NormalTok{\}}
\end{Highlighting}
\end{Shaded}

dove il primo metodo recupera la lista di tutti i prodotti, il secondo
di tutti quelli di genere maschile e il terzo di tutti quelli di genere
femminile. L'implementazione specifica potrebbe essere anche diversa,
con, ad esempio, i due metodi per i generi che chiamano un unico metodo
in cui il genere è un parametro, ma questo non è rilevante per chi li
utilizzerà, in quanto quello che conta è che, chiamandoli, riceverà in
cambio un vettore contenente i capi di abbigliamento desiderati.

\hypertarget{la-classe-prodottocontroller}{%
\subsection{La classe
ProdottoController}\label{la-classe-prodottocontroller}}

Il \emph{controller} si occuperà invece di far interagire la vista con
il modello, prendendo gli input forniti dalla vista, generalmente sotto
forma di parametri GET o POST, elaborando l'azione o le azioni opportune
per recuperare i dati di interesse dal database, e infine iniettando il
risultato nella vista opportuna.

Anche in questo caso un'implementazione ragionevole potrebbe essere la
seguente:

\begin{Shaded}
\begin{Highlighting}[]
\KeywordTok{class} \ConstantTok{P}\NormalTok{rodottoController\{}

    \KeywordTok{private} \ConstantTok{E}\NormalTok{ngine }\VariableTok{$templateEngine}\OtherTok{;}

    \KeywordTok{public} \KeywordTok{function} \BuiltInTok{\_\_construct}\NormalTok{(}\ConstantTok{E}\NormalTok{ngine }\VariableTok{$templateEngine}\NormalTok{)\{}
        \VariableTok{$this}\NormalTok{{-}\textgreater{}templateEngine }\OperatorTok{=} \VariableTok{$templateEngine}\OtherTok{;}
\NormalTok{    \}}

    \KeywordTok{public} \KeywordTok{function}\NormalTok{ listAll()\{}
        \ControlFlowTok{return} \VariableTok{$this}\NormalTok{{-}\textgreater{}listAllByGenre(}\StringTok{\textquotesingle{}All\textquotesingle{}}\NormalTok{)}\OtherTok{;}
\NormalTok{    \}}

    \KeywordTok{public} \KeywordTok{function}\NormalTok{ listAllByGenre(}\VariableTok{$genere}\NormalTok{)\{}
        \ControlFlowTok{if}\NormalTok{ (}\VariableTok{$genere} \OperatorTok{===} \StringTok{\textquotesingle{}All\textquotesingle{}}\NormalTok{)}
            \VariableTok{$prodotti} \OperatorTok{=} \ConstantTok{P}\NormalTok{rodottoRepository::listAll()}\OtherTok{;}
        \ControlFlowTok{else} \ControlFlowTok{if}\NormalTok{ (}\VariableTok{$genere} \OperatorTok{===} \StringTok{\textquotesingle{}Uomo\textquotesingle{}}\NormalTok{)}
            \VariableTok{$prodotti} \OperatorTok{=} \ConstantTok{P}\NormalTok{rodottoRepository::listAllMale()}\OtherTok{;}
        \ControlFlowTok{else}
            \VariableTok{$prodotti} \OperatorTok{=} \ConstantTok{P}\NormalTok{rodottoRepository::listAllFemale()}\OtherTok{;}
        \ControlFlowTok{return} \VariableTok{$this}\NormalTok{{-}\textgreater{}templateEngine{-}\textgreater{}render(}\StringTok{\textquotesingle{}negozio\textquotesingle{}}\OtherTok{,}
\NormalTok{            [}
                \StringTok{\textquotesingle{}prodotti\textquotesingle{}}\NormalTok{ =\textgreater{} }\VariableTok{$prodotti}\OtherTok{,}
                \StringTok{\textquotesingle{}genere\textquotesingle{}}\NormalTok{ =\textgreater{} }\VariableTok{$genere}
\NormalTok{            ]}
\NormalTok{        )}\OtherTok{;}

\NormalTok{    \}}
\NormalTok{\}}
\end{Highlighting}
\end{Shaded}

A differenza del modello, in questo caso si è preferito avere una classe
di cui istanziare degli oggetti e che avesse come unico attributo un
riferimento al \texttt{template}, sotto forma di un oggetto di tipo
\texttt{Engine}, visto in precedenza e che rappresenta il modo in cui
Plates gestisce la generazione di template.

Creando quindi un oggetto di tipo \texttt{ProdottoController}, nel
costruttore si passerà l'oggetto per la gestione del template, in modo
che nel codice dei vari metodi il controller possa utilizzarlo per il
render del codice HTML. Guardando ad esempio il metodo
\texttt{listAllByGenre}, il cui scopo è quello di restituire la lista
dei prodotti in base al genere desiderato, che può essere \emph{All},
\emph{Donna} o \emph{Uomo}, si vede che il codice si interfaccia con il
modello per chiedere l'elenco desiderato in base al genere, e,
successivamente, utilizza i dati recuperati per passarli al template che
si preoccuperà di farne la resa in HTML, ritornando proprio l'HTML
generato.

In questo esempio il template, \texttt{negozio.tpl}, avrà una struttura
di questo tipo:

\begin{Shaded}
\begin{Highlighting}[]
\DataTypeTok{\textless{}!doctype }\NormalTok{html}\DataTypeTok{\textgreater{}}
\DataTypeTok{\textless{}}\KeywordTok{html} \ErrorTok{lang}\OtherTok{=}\StringTok{"it"}\DataTypeTok{\textgreater{}}
\DataTypeTok{\textless{}}\KeywordTok{head}\DataTypeTok{\textgreater{}}
    \DataTypeTok{\textless{}}\KeywordTok{meta} \ErrorTok{charset}\OtherTok{=}\StringTok{"UTF{-}8"}\DataTypeTok{\textgreater{}}
    \DataTypeTok{\textless{}}\KeywordTok{title}\DataTypeTok{\textgreater{}}\NormalTok{Esempio negozio con pattern MVC}\DataTypeTok{\textless{}/}\KeywordTok{title}\DataTypeTok{\textgreater{}}
\DataTypeTok{\textless{}/}\KeywordTok{head}\DataTypeTok{\textgreater{}}
\DataTypeTok{\textless{}}\KeywordTok{body}\DataTypeTok{\textgreater{}}
    \DataTypeTok{\textless{}}\KeywordTok{h1}\DataTypeTok{\textgreater{}}\NormalTok{Esempio negozio con pattern MVC}\DataTypeTok{\textless{}/}\KeywordTok{h1}\DataTypeTok{\textgreater{}}
    \DataTypeTok{\textless{}}\KeywordTok{p}\DataTypeTok{\textgreater{}\textless{}}\KeywordTok{a} \ErrorTok{href}\OtherTok{=}\StringTok{"}\ErrorTok{\textless{}}\StringTok{?=$base\_path?\textgreater{}/negozio/genere/All"}\DataTypeTok{\textgreater{}}\NormalTok{Tutti i capi}\DataTypeTok{\textless{}/}\KeywordTok{a}\DataTypeTok{\textgreater{}\textless{}/}\KeywordTok{p}\DataTypeTok{\textgreater{}}
    \DataTypeTok{\textless{}}\KeywordTok{p}\DataTypeTok{\textgreater{}\textless{}}\KeywordTok{a} \ErrorTok{href}\OtherTok{=}\StringTok{"}\ErrorTok{\textless{}}\StringTok{?=$base\_path?\textgreater{}/negozio/genere/Donna"}\DataTypeTok{\textgreater{}}\NormalTok{Abbigliamento femminile}\DataTypeTok{\textless{}/}\KeywordTok{a}\DataTypeTok{\textgreater{}\textless{}/}\KeywordTok{p}\DataTypeTok{\textgreater{}}
    \DataTypeTok{\textless{}}\KeywordTok{p}\DataTypeTok{\textgreater{}\textless{}}\KeywordTok{a} \ErrorTok{href}\OtherTok{=}\StringTok{"}\ErrorTok{\textless{}}\StringTok{?=$base\_path?\textgreater{}/negozio/genere/Uomo"}\DataTypeTok{\textgreater{}}\NormalTok{Abbigliamento maschile}\DataTypeTok{\textless{}/}\KeywordTok{a}\DataTypeTok{\textgreater{}\textless{}/}\KeywordTok{p}\DataTypeTok{\textgreater{}}
    \DataTypeTok{\textless{}}\KeywordTok{h2}\DataTypeTok{\textgreater{}}\NormalTok{Lista dei prodotti: }\KeywordTok{\textless{}?}\NormalTok{=$genere}\KeywordTok{?\textgreater{}}\DataTypeTok{\textless{}/}\KeywordTok{h2}\DataTypeTok{\textgreater{}}
    \DataTypeTok{\textless{}}\KeywordTok{ul}\DataTypeTok{\textgreater{}}
        \KeywordTok{\textless{}?php}\NormalTok{ foreach ($prodotti as $prodotto): }\KeywordTok{?\textgreater{}}
            \DataTypeTok{\textless{}}\KeywordTok{li}\DataTypeTok{\textgreater{}}\KeywordTok{\textless{}?}\NormalTok{=$prodotto[\textquotesingle{}nome\textquotesingle{}]}\KeywordTok{?\textgreater{}}\NormalTok{: }\DataTypeTok{\textless{}}\KeywordTok{i}\DataTypeTok{\textgreater{}}\KeywordTok{\textless{}?}\NormalTok{=$prodotto[\textquotesingle{}descrizione\textquotesingle{}]}\KeywordTok{?\textgreater{}}\DataTypeTok{\textless{}/}\KeywordTok{i}\DataTypeTok{\textgreater{}\textless{}/}\KeywordTok{li}\DataTypeTok{\textgreater{}}
        \KeywordTok{\textless{}?php}\NormalTok{ endforeach;}\KeywordTok{?\textgreater{}}
    \DataTypeTok{\textless{}/}\KeywordTok{ul}\DataTypeTok{\textgreater{}}
\DataTypeTok{\textless{}/}\KeywordTok{body}\DataTypeTok{\textgreater{}}
\DataTypeTok{\textless{}/}\KeywordTok{html}\DataTypeTok{\textgreater{}}
\end{Highlighting}
\end{Shaded}

e la resa, con la rotta \texttt{GET\ /negozio/genere/Donna} potrebbe
essere la seguente pagina:

\begin{figure}
\centering
\includegraphics{esempio.png}
\caption{Esempio della pagina dell'elenco dei prodotti}
\end{figure}

\hypertarget{riassunto}{%
\section{Riassunto}\label{riassunto}}

Riassumendo quanto visto finora:

\begin{itemize}
\tightlist
\item
  l'applicazione ha un file che funge da \emph{front-controller}, cioè
  un file a cui arrivano tutte le richieste, indipendentemente dall'URL,
  che si preoccuperà di dirottarle al controller opportuno. Questo file
  sarà sempre il file \texttt{index.html} e sarà all'interno della
  cartella \texttt{public}
\item
  il server web, in questo caso Apache, deve essere configurato per
  ridirigere tutte le richieste al \emph{front-controller}
\item
  le configurazioni del database o altro vanno inserite nel file
  \texttt{config.php} all'interno della cartella \texttt{conf}
\item
  la connessione al database viene gestita dalla classe
  \texttt{Connection}
\item
  il modello dei dati viene incapsulato nella classe, o nelle classi,
  che ``parlano'' direttamente con il database e restituiscono, tramite
  metodi statici, dei vettori contenenti le informazioni desiderate. Per
  convenzione i modelli verranno messi nella cartella \texttt{Model} e
  saranno chiamati con il nome che rappresenta i dati, ad esempio
  \emph{Prodotto}, seguito dalla parola \emph{Repository}
\item
  il controller, che si occupa dell'interazione con la vista e il
  modello, sarà una classe, o più classi, contenente un riferimento al
  template, che verrà referenziato da un parametro all'atto della
  creazione dell'oggetto, e i metodi per rappresentare le funzionalità
  dell'applicazione. Per convenzione i controller verranno messi nella
  cartella `Controller' e saranno chiamati con un nome che ne
  rappresenta la funzionalità, ad esempio \emph{Prodotto}, seguito dalla
  parola \emph{Controller}
\end{itemize}

La struttura delle cartelle sarà quindi la seguente:

\begin{figure}
\centering
\includegraphics{struttura.png}
\caption{Struttura delle cartelle di un'applicazione Slim organizzata
secondo il pattern MVC}
\end{figure}
